\documentclass{article}
\usepackage{amsmath,amssymb}
\begin{document}
\emph{Fixed-trace beta deformation of the exchangeable weight on the profile simplex.} Within the named coupled gaussian slot law and Haar frame measure of the preceding paragraph, and at fixed \(D\) and unit trace, this paragraph records what follows for the exchangeable profile weights with density \(\prod_{i}p_{i}^{a-1}\prod_{i<j}\lvertz p_{i}-p_{j}\rvert^{\beta}\) on the simplex, taken as a named input family with parameters \(a>0\) and \(\beta\ge 0\): the factor \(\prod_{i<j}\lvert p_{i}-p_{j}\rvert^{\beta}\) is posed here as a deformation, and no preceding computation yields it from the named slot or frame measures. Two endpoints of the preceding paragraph sit in the family at \(a=1\), \(\beta=0\): the flat Haar frame weight, and the slot pushforward at uniform profile. Writing \(E_{i}=\tau p_{i}\) with \(p_{D}=1-\sum_{i<D}p_{i}\), the Jacobian of \((E_{1},\dots,E_{D})\mapsto(\tau,p_{1},\dots,p_{D-1})\) is exactly \(\tau^{D-1}\), so the fixed-trace weight with exponential scale factor \(e^{-E_{i}/\sigma}\) and the deformation factor separates exactly into \(\sigma^{Da+\beta D(D-1)/2}\Gamma(Da+\beta D(D-1)/2)\,\tau^{Da+\beta D(D-1)/2-1}e^{-\tau/\sigma}\) times the profile weight above, and the profile marginal carries no scale parameter. On the chamber \(p_{1}>\dots>p_{D}>0\) the log-weight is strictly concave for \(a\ge 1\), so for \(a>1\) it has exactly one stationary profile there; and at any interior stationary profile the multiplier of the trace constraint is forced by the homogeneity count \(\sum_{i}p_{i}\partial_{i}L=D(a-1)+\beta D(D-1)/2\) to equal that constant, while for \(a=1\) the summed-gradient identity \(\sum_{i}\partial_{i}L=0\) is incompatible with that multiplier for every \(\beta>0\): the \(a=1\) members have no interior stationary profile at any deformation strength, and for \(D=2\) the density \((2t-1)^{\beta}\) is maximized only at the vertices. The family is continuous at \(\beta=0\): for \(a=1\) its normalizations tend to the flat values, \(1/(\beta+1)\to 1\) at \(D=2\) and \(1/2\) at \(D=3\), by dominated convergence from \(\prod_{i<j}\lvert p_{i}-p_{j}\rvert^{\beta}\le 1\). For \(D=2\) and \(a>1\) the stationary profile is known in closed form: the separation \(u=p_{1}-p_{2}\) is \(u=\sqrt{\beta/(\beta+2(a-1))}\) with log-weight curvature \(-(2a+\beta-2)^{2}/(a-1)\), so the departure of the stationary profile from the barycenter grows like \(\sqrt{\beta/(2(a-1))}\) at small \(\beta\), while at large \(\beta\) the gap to the vertex obeys \((1-u)\beta\to a-1\); the same square-root onset is found numerically at \(D=3\). For \(D=3\) and \(D=4\) at \(a=2\), \(\beta=1\) the stationary profiles are \((0.6466,0.2755,0.0780)\) and \((0.5477,0.2866,0.1286,0.0372)\), with stationarity residuals below \(10^{-28}\) and, at \(D=3\), the grid maximum of the weight within one grid step of the stationary profile. The deformed weight is not flat: at \(a=2\), \(\beta=1\) it vanishes at the barycenter, is positive at the stationary profile, and weights the largest entry with mean \(0.629\) against the flat value \(1/3\), checked on two independent numerical paths. Every member is relabelling-invariant, so its extremal profiles form full relabelling orbits and no single member of an orbit is picked out inside the family: the deformation displaces weight from the uniformity locus and creates, exactly for \(a>1\), a unique chamber profile up to relabelling, realizing the relabelling-symmetric alternative of the second preceding paragraph rather than supplying a selector; it supplies no state, its origin in the matrix integral is not established, it leaves the missing steps (i)-(iv) of the Missing step paragraph unchanged, and it closes nothing in Problem~G.
\end{document}
